\begin{helper}[FiberAndDegreeSameColorBaseCodegreeSharpFailureNegligible]
\label{helper:FiberAndDegreeSameColorBaseCodegreeSharpFailureNegligible}
Let \(\beta=1/2\), let
\[
m(n)=\noderef{TwoBiteNaturalReducedVertexCount}(n),
\qquad
p(n)=\noderef{TwoBiteEdgeProbability}(\beta,n),
\]
and take probabilities with respect to
\(\noderef{TwoBiteEventProbability}(n,m(n),p(n))\).  The probability that
the sharp same-color base-codegree event fails is
\(\noderef{NegligibleProbability}\).  Explicitly, the negligible event is the
failure of
\[
\begin{aligned}
&\forall r\ne s,\quad
|\{t:\noderef{TwoBiteRedGraph}(\omega).Adj(r,t)\land
\noderef{TwoBiteRedGraph}(\omega).Adj(s,t)\}|\le \log n,\\
&\forall b\ne c,\quad
|\{t:\noderef{TwoBiteBlueGraph}(\omega).Adj(b,t)\land
\noderef{TwoBiteBlueGraph}(\omega).Adj(c,t)\}|\le \log n .
\end{aligned}
\]
\end{helper}

\begin{proof}
Write \(M=m(n)\), \(p=p(n)\), and \(L=\log n\).  It is enough to discard
finitely many \(n\).  From \(\noderef{TwoBiteNaturalReducedVertexCount}\),
\(\noderef{TwoBiteReducedVertexCount}\), and \(\noderef{TwoBiteStretch}\), we
have \(M=\lfloor n/L^2\rfloor\) for all large \(n\), and hence \(M\le n\).
By \(\noderef{OppositeProjectionEdgeProbBounds}\), also \(0\le p\le1\)
eventually.  We also use \(\noderef{TwoBiteNaturalTotalMassOneEventually}\):
eventually
\[
n\le M^2
\]
and the full \(\noderef{TwoBiteEventProbability}\) mass of the two-bite
sample space is \(1\).  The inequality \(n\le M^2\) is the support condition
needed below for the injection coordinate.

Fix distinct red vertices \(r,s\), and let \(E_{r,s}\) be the event that their
red base common-neighbor count is larger than \(L\).  A sample
\(\omega\in\noderef{TwoBiteSample}(n,M,p)\) is a triple
\((G_R,G_B,\pi)\).  By \(\noderef{TwoBiteEventProbability}\) and
\(\noderef{TwoBiteSampleWeight}\),
\[
\Pr(E_{r,s})
=
\sum_{G_R,G_B,\pi}
\mathbf 1_{E_{r,s}}(G_R)\,
\noderef{GnpGraphWeight}(p,G_R)\,
\noderef{GnpGraphWeight}(p,G_B)\,
\noderef{UniformInjectionWeight}(n,M).
\tag{1}
\]
The indicator depends only on \(G_R\).  Because \(n\le M^2\), injections
\(\operatorname{Fin}(n)\hookrightarrow\operatorname{Fin}(M)\times
\operatorname{Fin}(M)\) exist, so \(\noderef{UniformInjectionWeight}\) is the
reciprocal of the number of such injections.  Hence
\[
\sum_\pi \noderef{UniformInjectionWeight}(n,M)=1.
\tag{2}
\]
Also
\[
\sum_{G_B}\noderef{GnpGraphWeight}(p,G_B)=1
\tag{3}
\]
by \(\noderef{GnpGraphWeightSumOne}\).  Applying (2) and (3) to (1) gives the
red marginal formula
\[
\Pr(E_{r,s})
=
\sum_{G_R:E_{r,s}(G_R)}
\noderef{GnpGraphWeight}(p,G_R).
\tag{4}
\]
Thus the fixed red pair has exactly the usual \(G(M,p)\) product law.  The
same finite-sum calculation, with the two base graph coordinates interchanged,
gives the identical formula for a fixed blue pair.

Under the \(G(M,p)\) product law, the common-neighbor count of a fixed pair is
stochastically dominated by a binomial variable with \(M\) trials and success
probability \(p^2\): for any proposed common-neighbor set \(S\) of size \(T\),
requiring all vertices of \(S\) to be common neighbors requires the two
independent edge events from each \(t\in S\), and therefore has weight at most
\((p^2)^T\).  Let \(T=\lceil L\rceil\).  If the count is \(>L\), then it is at
least \(T\), so the union bound over \(T\)-subsets gives
\[
\Pr(E_{r,s})
\le \binom MT(p^2)^T.
\tag{5}
\]
Since \(\beta=1/2\), \(\noderef{TwoBiteEdgeProbability}\) gives
\[
p^2=\frac14\,\frac{L}{n}.
\]
With \(M\le n/L^2\), the binomial mean \(\mu=Mp^2\) satisfies
\[
\mu\le \frac{1}{4L}.
\tag{6}
\]
Using \(\noderef{OppositeProjectionChooseExponentialBound}\),
\[
\binom MT(p^2)^T
\le \left(\frac{eMp^2}{T}\right)^T
= \left(\frac{e\mu}{T}\right)^T.
\tag{7}
\]
Because \(T\ge L\), (6) implies
\[
\left(\frac{e\mu}{T}\right)^T
\le
\left(\frac{e}{4L^2}\right)^L
\le \exp\!\left(-\frac12 L\log L\right)
\tag{8}
\]
for all sufficiently large \(n\).

There are fewer than \(M^2\le n^2=\exp(2L)\) ordered red pairs and the same
number of ordered blue pairs.  Applying the union bound over all same-color
pairs and using (8) gives the bound
\[
2\exp(2L)\exp\!\left(-\frac12 L\log L\right).
\]
This tends to \(0\), because \(L\log L\) dominates \(L\).  Therefore the
probability that the sharp same-color base-codegree event fails tends to
\(0\), which is exactly the asserted \(\noderef{NegligibleProbability}\).
\end{proof}
